\section{What the Literature Cannot Yet Answer: A Research Agenda}\label{sec:agenda}

The framework's weakest links are questions that matter most for unattended systems. We frame six as studies. Each should meet the standard of Table~\ref{tab:reporting}, including at least three runs per arm, budget-matched arms, held-out tasks and a grader the systems under test cannot reach.

\textbf{Does a builder paired with an independent tester beat a budget-matched solo agent on applications built from scratch?} The arms are a solo agent, a solo agent instructed to test its own application, and a builder plus a separate tester who uses the running application, at equal token or dollar budgets. The metric is the share of hidden workflow checks passed, with pass\textasciicircum k and cost per passed task. The closest controlled studies are mixed: in end-to-end development, a developer followed by a tester led a single agent on average but not with one of two models \citep{T12}, and a full team was not significantly better than a solo agent \citep{T18}. The browser-testing result is a cross-model correlation \citep{C19}, and the one long-running comparison with an independent evaluator in our corpus, an Anthropic engineering post, is a single run, judged by its author, at more than 20 times the solo cost \citep{T28}.

\textbf{For runs that last hours, are separate sessions better than one long session?} The arms are one session with automatic compaction and a chain of fresh sessions that hand off through files and version control, at equal budget. The metrics are tasks completed, interventions per task, and work lost or redone after a reset. Errors left in context degrade later turns \citep{F3}, which favors fresh sessions, but handoffs lose information \citep{T18,T22}, which favors one session, and notes carried between runs gave a small gain \citep{H20}.

\textbf{Does an honest ``stuck'' exit reduce gaming without reducing completion?} The arms cross an exit option, present or absent, with solvable and deliberately impossible tasks. The metrics are cheating on impossible tasks, and completion and false flags on solvable ones. In the one study, the exit cut cheating sharply for OpenAI models but much less for Claude Opus~4.1, and its authors report a relatively small loss on the original benchmark \citep{V3}; whether it holds on hard but solvable tasks is open.

\textbf{How much does pre-registration reduce forking paths in agent-run analysis?} The arms are free analysis, a plan committed before the agent sees the data, and a committed plan plus an independent re-run, on tasks with known answers. The metrics are the spread of estimates and verdicts across runs, the distance from the known answer, and how often verdicts flip under confirmatory prompting. Current studies mainly measure how far agents can be steered \citep{R2,R5}; in one, an AI auditor narrowed the spread of conclusions without eliminating it \citep{R1}. Evidence that pre-registration helps is observational and comes from human researchers \citep{R13,R19}, and the one agent result is a retrospective counterfactual \citep{R6}.

\textbf{When is hardening a grader worth its cost?} The arms are a grader before and after adversarial hardening, on the same tasks and agents. The metrics are the hackable share, false rejections of correct reference solutions, changes in agent rankings, and the dollars and hours hardening takes. Existing studies show that such loops cut exploitability \citep{V27,T24}; a decision to harden needs that gain weighed against these costs.

\textbf{How fast do reliability and harness value change across model versions?} The design holds harness, tasks and grader fixed across successive model versions, switching harness components on and off for each. The metrics are pass\textasciicircum k, cost per passed task and the gain from each component. Results are version-specific (Section~\ref{sec:scope}): one harness component helped an older model more than its successor \citep{H20}, and capability gains have brought only small improvements in reliability \citep{F2}.

A seventh task underlies the rest: coding the corpus against Table~\ref{tab:reporting}, a count this review did not attempt (Section~\ref{sec:measurement}).
